% ---
% Subsección nueva en TT2 (2026-10-18), al final del Marco Teórico para no renumerar lo aprobado.
% Reúne los métodos clásicos y de kernel que usan la Metodología y los Resultados: la radiómica (M3),
% PCA y Kernel PCA (C2 y C3, D-016, D-027), la SVM y el QSVM (X4, D-038) y el kernel proyectado (X6).
% ---

\subsection{Representaciones clásicas y modelos de kernel}
\label{sec:representaciones_clasicas}

Las secciones anteriores describen la parte cuántica del sistema. Esta reúne los métodos clásicos que la rodean: la extracción radiómica, que produce el vector de características de cada lesión; las dos transformaciones clásicas con las que se comparan las cuánticas, el análisis de componentes principales y su versión con kernel, y los modelos de kernel que se usan en los experimentos complementarios.
